%!TEX program = xelatex
\documentclass[11pt]{ctexart}
% Shared packages / theorem environments for standalone topic articles.
% Body content must live in each topic's .tex (do not \input chapters/).

\usepackage[a4paper,margin=2.35cm]{geometry}
\usepackage{amsmath,amssymb,amsthm,mathtools}
\usepackage{bm}
\usepackage{enumitem}
\usepackage{booktabs}
\usepackage{longtable}
\usepackage{array}
\usepackage{xcolor}
\usepackage{tikz}
\usetikzlibrary{trees,positioning,shapes.geometric,calc}
\usepackage{xurl}
\usepackage{hyperref}
\usepackage{csquotes}
\usepackage[backend=biber,style=authoryear,sorting=nyt,maxcitenames=2,maxbibnames=6]{biblatex}

\addbibresource{refs.bib}

\hypersetup{
  colorlinks=true,
  linkcolor=blue!50!black,
  citecolor=blue!50!black,
  urlcolor=blue!60!black
}

\setlist[itemize]{topsep=3pt,itemsep=2pt,parsep=0pt}
\setlist[enumerate]{topsep=3pt,itemsep=2pt,parsep=0pt}
\setcounter{tocdepth}{2}
\setlength{\emergencystretch}{3em}
\raggedbottom
\newcolumntype{L}[1]{>{\raggedright\arraybackslash}p{#1}}

% --- standard amsthm environments (section-numbered) ---
\theoremstyle{plain}
\newtheorem{theorem}{定理}[section]
\newtheorem{lemma}[theorem]{引理}
\newtheorem{proposition}[theorem]{命题}
\newtheorem{corollary}[theorem]{推论}

\theoremstyle{definition}
\newtheorem{definition}{定义}[section]
\newtheorem{example}[definition]{例}
\newtheorem{teaching}[definition]{从零教学}

\theoremstyle{remark}
\newtheorem{remark}{注}[section]
\newtheorem{heuristic}[remark]{启发式解释}
\newtheorem{engineering}[remark]{工程惯例}
\newtheorem{examnote}[remark]{考试提醒}
\newtheorem{checkpoint}[remark]{考点}
\newtheorem{proofsketch}[remark]{证明骨架}

% Keep examproblem available if a topic ever needs it
\newtheoremstyle{examstyle}
  {6pt}{6pt}{\normalfont}{}{\bfseries}{.}{0.5em}{}
\theoremstyle{examstyle}
\newtheorem{examproblem}{题}[section]

% --- math macros ---
\newcommand{\R}{\mathbb{R}}
\newcommand{\E}{\mathbb{E}}
\newcommand{\Pbb}{\mathbb{P}}
\newcommand{\N}{\mathcal{N}}
\newcommand{\Hcal}{\mathcal{H}}
\newcommand{\Dcal}{\mathcal{D}}
\newcommand{\Lcal}{\mathcal{L}}
\newcommand{\argmin}{\operatorname*{arg\,min}}
\newcommand{\argmax}{\operatorname*{arg\,max}}
\newcommand{\prox}{\operatorname{prox}}
\newcommand{\Var}{\operatorname{Var}}
\newcommand{\KL}{\operatorname{KL}}
\newcommand{\softmax}{\operatorname{softmax}}
\newcommand{\Dir}{\operatorname{Dir}}
\newcommand{\Cat}{\operatorname{Cat}}
\newcommand{\Bern}{\operatorname{Bern}}
\newcommand{\tr}{\operatorname{tr}}

\newcommand{\source}[1]{\par\smallskip\noindent{\small\textit{来源定位：#1}}\par\smallskip}
\newcommand{\solution}{\par\smallskip\noindent\textbf{解答。}\ }
\newcommand{\officialenglish}[1]{\par\smallskip\noindent\textbf{Official English prompt.}\ \begin{quote}\small #1\end{quote}}
\newcommand{\chinesetranslation}[1]{\par\smallskip\noindent\textbf{中文翻译。}\ #1\par\smallskip}
\newcommand{\concepttarget}[1]{\phantomsection\label{#1}\hypertarget{#1}{}}
\newcommand{\conceptref}[2]{\hyperref[#1]{#2}}


\title{\bfseries 求真书院 AI 博士生资格考试\\模块一：考试大纲与参考书目 (Syllabus \& Textbooks)}
\author{清华大学求真书院 AI 方向博士生资格考试备考资料}
\date{\today}

\begin{document}

\maketitle
\tableofcontents

\section{说明与使用方式}

本资料面向清华大学求真书院人工智能方向博士生资格考试（QE-AI），按公开试卷、考试大纲和指定教材整理；具体出处见文末参考文献。

题目按题号、主题和关键条件组织，不逐字复制原卷全文；核对题意时应以校方发布版本为准。

\subsection{如何复习}

四个方向各 33 分，考试通常要求四选三，另有 1 分身份信息分。复习时建议分三轮：

\begin{enumerate}
  \item 第一轮读本讲义与理论预备材料，建立四门核心课的共同符号。
  \item 第二轮按年份刷各年份真题解答，先自己写解，再对照推导。
  \item 第三轮按考前 checklist 反向查漏，重点补 PAC/VC、prox-SGD、扩散模型、policy gradient、Transformer 和 RAG/MLN。
\end{enumerate}

\subsection{解答风格}

每道题按以下逻辑写：

\begin{itemize}
  \item 先说明题意定位，避免和原卷 PDF 的具体排版绑定；
  \item 再给核心公式或构造；
  \item 需要证明时给可独立复现的推导；
  \item 对有歧义的题，明确“严格结论”和“考试语境下常见答案”。
\end{itemize}

本文不是替代教材。大纲指定的机器学习理论主线可读 \textcite{shalev2014uml,mohri2018foundations}，深度学习和生成模型可读 \textcite{goodfellow2016deep,bishop2023deep,bishop2006prml}，优化可读 \textcite{nesterov2018convex,beck2017first}，NLP 可读 \textcite{jurafsky2026slp,goldberg2017nnnlp}。


\section{考试大纲与考点地图}
\label{sec:syllabus}

求真书院 2026 年 AI 方向博士生资格考试大纲列出四门核心课程：Machine Learning Theory、Advanced Deep Learning、Optimization Methods for AI、Natural Language Understanding。这四门课也是三套公开试卷的四个 Section。以下按官方大纲的条目顺序给出复习地图；每项的定义、推导和题型训练在理论预备与各年份真题解答中展开 \parencite{qzc-ai-syllabus-2026}。

\newcommand{\MLTextbooks}{
\subsubsection{Understanding Machine Learning: From Theory to Algorithms}

\textcite{shalev2014uml} 是本课程从学习问题形式化到可学习性证明的主线教材。尤其应把第 2--7 章连成 ERM--PAC--uniform convergence--VC--SRM 的证明链。

\begingroup
\small
\begin{longtable}{L{0.22\textwidth}L{0.50\textwidth}L{0.18\textwidth}}
\toprule
原书章节群 & 讲什么 & 对应考纲/考点 \\
\midrule
\endfirsthead
\toprule
原书章节群 & 讲什么 & 对应考纲/考点 \\
\midrule
\endhead
第 1--2 章：学习问题与 ERM & 统计学习框架、经验风险、过拟合、归纳偏置、有限假设类。 & 监督学习、经验/真实风险。 \\
第 3--4 章：PAC 与 uniform convergence & realizable/agnostic PAC、样本复杂度，以及有限类的泛化保证。 & PAC 定义、Hoeffding/union bound 证明套路。 \\
第 5--7 章：No-Free-Lunch、VC 与模型选择 & 误差分解、VC 维、Sauer 引理、基本学习定理、SRM 与 MDL。 & VC/Sauer、偏差--方差、模型选择。 \\
第 8--14 章：可计算学习与算法例子 & ERM 的计算复杂度、半空间、线性回归、正则化、核方法、决策树和最近邻。 & 回归、SVM/核、树、\(k\)-NN。 \\
\bottomrule
\end{longtable}
\endgroup

\subsubsection{Foundations of Machine Learning, 2nd ed.}

\textcite{mohri2018foundations} 提供更偏理论和算法的补充视角；对于 Rademacher complexity、margin、kernel 和 online learning，它比前书更适合作为证明细节的回查来源。

\begingroup
\small
\begin{longtable}{L{0.22\textwidth}L{0.50\textwidth}L{0.18\textwidth}}
\toprule
原书章节群 & 讲什么 & 对应考纲/考点 \\
\midrule
\endfirsthead
\toprule
原书章节群 & 讲什么 & 对应考纲/考点 \\
\midrule
\endhead
第 1--2 章：学习任务与 PAC & 学习阶段、泛化、有限假设集的可实现/不可实现保证、噪声与 Bayes error。 & PAC 框架与统计假设。 \\
第 3--4 章：复杂度与模型选择 & Rademacher complexity、growth function、VC 维下界、ERM/SRM、交叉验证、正则化和 surrogate loss。 & VC 泛化界、模型选择。 \\
第 5--6 章：margin 与核 & 线性分类、support vector、对偶问题、margin theory、RKHS、representer theorem、sequence kernel。 & SVM、kernel trick、正定性。 \\
第 7 章以后：boosting 与在线学习 & AdaBoost、\(\ell_1\) 正则、专家建议、mistake/regret 类保证及后续学习算法专题。 & 了解泛化保证如何连接算法。 \\
\bottomrule
\end{longtable}
\endgroup

}

\newcommand{\DeepTextbooks}{
\subsubsection{Deep Learning}

\textcite{goodfellow2016deep} 的三部分结构依次给出数学与 ML 基础、深度网络训练、表示与生成模型，是神经网络、优化和生成模型的统一参考。

\begingroup
\small
\begin{longtable}{L{0.22\textwidth}L{0.50\textwidth}L{0.18\textwidth}}
\toprule
原书章节群 & 讲什么 & 对应考纲/考点 \\
\midrule
\endfirsthead
\toprule
原书章节群 & 讲什么 & 对应考纲/考点 \\
\midrule
\endhead
第 1--5 章：基础 & 线性代数、概率/信息论、数值计算、容量与泛化、MLE/MAP、SGD。 & 损失、梯度、泛化、概率建模。 \\
第 6--8 章：深网络与训练 & 前馈网络、反向传播、正则化、优化、初始化、batch normalization。 & MLP、训练方法、稳定性。 \\
第 9--12 章：结构与实践 & CNN、序列模型、practical methodology、应用。 & 视觉特征、序列/注意力的前置知识。 \\
第 13--20 章：表示和生成 & 自编码器、结构化概率模型、Monte Carlo、近似推断、深度生成模型。 & 表征学习、VAE/GAN/flow 的理论背景。 \\
\bottomrule
\end{longtable}
\endgroup

\subsubsection{Deep Learning: Foundations and Concepts}

\textcite{bishop2023deep} 以概率建模贯穿回归、分类、深网络、视觉、Transformer 与生成模型；其第 17--20 章尤其适合对照大纲中的视觉生成模型。

\begingroup
\small
\begin{longtable}{L{0.22\textwidth}L{0.50\textwidth}L{0.18\textwidth}}
\toprule
原书章节群 & 讲什么 & 对应考纲/考点 \\
\midrule
\endfirsthead
\toprule
原书章节群 & 讲什么 & 对应考纲/考点 \\
\midrule
\endhead
第 1--5 章：概率与单层网络 & 深度学习概览、概率/信息论、分布、线性回归、判别/生成分类。 & 似然、交叉熵、回归/分类基础。 \\
第 6--9 章：深网络训练 & 多层网络、gradient descent、backpropagation、正则化与 residual connection。 & 网络训练、泛化与归纳偏置。 \\
第 10--13 章：视觉、序列和 Transformer & CNN、检测/分割/风格迁移、图模型和序列、attention/Transformer、图网络。 & 视觉任务、Transformer 实现。 \\
第 14--20 章：采样和生成 & MCMC、离散/连续隐变量、GAN、normalizing flow、autoencoder/VAE、diffusion/score matching。 & 生成模型的目标、ELBO、采样。 \\
\bottomrule
\end{longtable}
\endgroup

\subsubsection{Pattern Recognition and Machine Learning}

\textcite{bishop2006prml} 的目录从概率分布、线性模型扩展到图模型与近似推断；它是把概率模型、优化目标和推断算法连起来的经典参考。

\begingroup
\small
\begin{longtable}{L{0.22\textwidth}L{0.50\textwidth}L{0.18\textwidth}}
\toprule
原书章节群 & 讲什么 & 对应考纲/考点 \\
\midrule
\endfirsthead
\toprule
原书章节群 & 讲什么 & 对应考纲/考点 \\
\midrule
\endhead
第 1--2 章：概率与决策 & 曲线拟合、Bayes 决策、信息论、指数族、非参数密度和最近邻。 & 概率建模、模型选择。 \\
第 3--4 章：线性回归和分类 & 最小二乘、bias--variance、Bayesian regression、logistic regression、Laplace approximation。 & 回归、逻辑回归、泛化。 \\
第 5--7 章：神经网络与核 & 反向传播、正则化、CNN、Gaussian process、SVM/RVM。 & MLP、核方法、margin。 \\
第 8--14 章：概率图模型与推断 & 图模型、EM、mixture、变分推断、sampling、序列模型、模型组合。 & 隐变量、EM、推断和生成建模。 \\
\bottomrule
\end{longtable}
\endgroup

\subsubsection{Computer Vision: Algorithms and Applications, 2nd ed.}

\textcite{szeliski2022vision} 是视觉任务的结构化目录：每一类任务都可按输入、几何/统计假设、输出以及评价指标复述。当前本库保留作者官方访问入口；表中目录按该版正式目录归并。

\begingroup
\small
\begin{longtable}{L{0.22\textwidth}L{0.50\textwidth}L{0.18\textwidth}}
\toprule
原书章节群 & 讲什么 & 对应考纲/考点 \\
\midrule
\endfirsthead
\toprule
原书章节群 & 讲什么 & 对应考纲/考点 \\
\midrule
\endhead
导论与成像 & 相机/图像形成、颜色、几何变换和视觉系统的基本假设。 & 视觉问题的输入、坐标和观测模型。 \\
特征、匹配与对齐 & 特征检测/描述、匹配、鲁棒估计、image alignment。 & 特征提取、配准、恢复。 \\
三维和运动 & 多视几何、structure from motion、dense motion/optical flow。 & 三维重构、光流估计。 \\
分割、识别和计算摄影 & 图像分割、目标/场景识别、深度视觉、computational photography。 & 识别、分割及评价指标。 \\
\bottomrule
\end{longtable}
\endgroup

}

\newcommand{\OptimizationTextbooks}{
\subsubsection{Lectures on Convex Optimization}

\textcite{nesterov2018convex} 强调凸优化的几何、复杂度和一阶方法；与真题相连时，应优先抓住光滑/强凸不等式、oracle 模型与收敛率之间的逻辑。

\begingroup
\small
\begin{longtable}{L{0.22\textwidth}L{0.50\textwidth}L{0.18\textwidth}}
\toprule
原书章节群 & 讲什么 & 对应考纲/考点 \\
\midrule
\endfirsthead
\toprule
原书章节群 & 讲什么 & 对应考纲/考点 \\
\midrule
\endhead
凸分析基础 & 凸集/凸函数、次梯度、分离与最优性条件。 & 凸性判别、次微分。 \\
复杂度与 oracle & 一阶 oracle、下界与不同函数类的可达到复杂度。 & 为什么要区分 smooth/strongly convex。 \\
一阶与加速方法 & gradient、accelerated gradient、mirror/prox 思想及收敛率。 & GD、动量/Nesterov 加速。 \\
约束与随机专题 & 约束处理、随机梯度和大规模优化的理论视角。 & SGD、随机/自适应优化的边界。 \\
\bottomrule
\end{longtable}
\endgroup

\subsubsection{First-Order Methods in Optimization}

\textcite{beck2017first} 的前六章给出本课程最常用的分析语言；第 8--15 章再把这些语言落实到 projected/subgradient、mirror、proximal、block、Frank--Wolfe 和 ADMM。

\begingroup
\small
\begin{longtable}{L{0.22\textwidth}L{0.50\textwidth}L{0.18\textwidth}}
\toprule
原书章节群 & 讲什么 & 对应考纲/考点 \\
\midrule
\endfirsthead
\toprule
原书章节群 & 讲什么 & 对应考纲/考点 \\
\midrule
\endhead
第 1--4 章：凸分析工具箱 & 向量空间、凸函数、subgradient、共轭函数和 Fenchel duality。 & 凸集/函数、次微分及性质。 \\
第 5--6 章：光滑性与 proximal & descent lemma、strong convexity、prox、Moreau envelope/decomposition。 & GD 收敛、复合优化。 \\
第 8--10 章：一阶核心算法 & projected/stochastic subgradient、mirror descent、proximal gradient、FISTA。 & SGD、随机坐标、动量加速。 \\
第 11--15 章：分块与分裂 & block proximal gradient、dual proximal、conditional gradient、alternating minimization、ADMM。 & 大规模优化的算法比较。 \\
\bottomrule
\end{longtable}
\endgroup

}

\newcommand{\NLPTextbooks}{
\subsubsection{Speech and Language Processing}

\textcite{jurafsky2026slp} 覆盖从统计 NLP 到现代大模型系统；针对本考纲，优先顺着“语言模型--embedding--Transformer--对齐/检索”的主线阅读。

\begingroup
\small
\begin{longtable}{L{0.22\textwidth}L{0.50\textwidth}L{0.18\textwidth}}
\toprule
原书章节群 & 讲什么 & 对应考纲/考点 \\
\midrule
\endfirsthead
\toprule
原书章节群 & 讲什么 & 对应考纲/考点 \\
\midrule
\endhead
第 1--5 章：统计 NLP 基础 & 文本规范化、edit distance、\(n\)-gram、smoothing、Naive Bayes、logistic regression。 & 统计建模假设、MLE/smoothing。 \\
第 6--10 章：向量和序列网络 & word vectors、神经语言模型、序列标注、RNN/LSTM 与 encoder--decoder。 & Word2vec、表示学习。 \\
Transformer 与大模型章节群 & attention、Transformer LM、预训练、instruction/alignment、评估和风险。 & self-attention、预训练对齐。 \\
语义、检索、对话与语音章节群 & 信息抽取、问答、dialogue、speech recognition/TTS 及应用系统。 & embedding 知识库、RAG 的任务背景。 \\
\bottomrule
\end{longtable}
\endgroup

\subsubsection{Neural Network Methods in Natural Language Processing}

\textcite{goldberg2017nnnlp} 是神经 NLP 的紧凑实现导向教材：从线性分类和计算图出发，逐步推到 embedding、CNN、RNN、attention 和条件生成。

\begingroup
\small
\begin{longtable}{L{0.22\textwidth}L{0.50\textwidth}L{0.18\textwidth}}
\toprule
原书章节群 & 讲什么 & 对应考纲/考点 \\
\midrule
\endfirsthead
\toprule
原书章节群 & 讲什么 & 对应考纲/考点 \\
\midrule
\endhead
第 1--2 章：任务和线性模型 & NLP 的离散结构、监督分类、one-hot/dense representation、损失、正则和 SGD。 & 统计假设、训练目标。 \\
第 3--5 章：MLP 与训练 & 非线性、前馈网络、embedding layer、computation graph、backprop、初始化和梯度问题。 & 神经网络实现、反向传播。 \\
第 6--10 章：文本表示与序列 & 特征到输入、language model、pretrained embeddings、CNN n-gram detector、RNN。 & Word2vec、序列建模。 \\
第 11--21 章：条件生成与结构 & encoder--decoder、attention、recursive network、multi-task/semi-supervised learning。 & attention、迁移和系统设计。 \\
\bottomrule
\end{longtable}
\endgroup
}


\subsection{机器学习理论}

官方大纲要求覆盖八组主题：

\begin{enumerate}
  \item 监督学习与 PAC 学习理论；
  \item VC 维、Sauer 引理，以及样本复杂度与 VC 维的关系；
  \item 模型选择与偏差--方差权衡；
  \item 线性回归与逻辑回归；
  \item 决策树；
  \item 最近邻算法；
  \item 支持向量机与核方法；
  \item 压缩感知。
\end{enumerate}

考纲指定教材为 \textcite{shalev2014uml} 和 \textcite{mohri2018foundations}；两条引用均链接至文末的标准书目信息与原书官网。

机器学习理论卷最稳定的考法不是背模型名称，而是证明学习保证。

必会对象如下：

\begin{itemize}
  \item 经验风险 \(L_S(h)\) 与真实风险 \(L_{\Dcal,f}(h)\)；
  \item VC 维、增长函数 \(\tau_{\Hcal}(m)\)、Sauer 引理；
  \item 可实现 PAC 与不可实现/agnostic PAC；
  \item 模型选择、偏差--方差、线性/逻辑回归、树模型与 \(k\)-NN；
  \item 核函数的正定性、Hilbert 空间特征映射、SVM margin 与 hinge loss；
  \item 稀疏测量、\(\ell_1\) 松弛与压缩感知的恢复条件。
\end{itemize}

\MLTextbooks

\subsection{深度学习和强化学习}

官方大纲把本部分分为五组：

\begin{enumerate}
  \item 神经网络框架、训练方法、逼近与泛化分析；
  \item 视觉任务中的特征提取、图像恢复、三维重构、光流估计、识别和分割；
  \item 无监督学习中的自编码器、对比学习、隐式正则化、迁移学习和元学习；
  \item 视觉生成模型中的 GAN、VAE、扩散模型、流模型及表达力分析；
  \item 强化学习中的贝尔曼方程、Q-learning 与 policy gradient。
\end{enumerate}

指定教材依次包括 \textcite{goodfellow2016deep,bishop2023deep,bishop2006prml,szeliski2022vision}。其中视觉任务应按输入、输出和评价指标组织，而生成模型与强化学习应能写出目标函数或递推式。

三套卷的重点可按上述条目定位：

\begin{itemize}
  \item MLP、CNN、Transformer 的结构、参数量、反向传播、归一化与泛化；
  \item 表征学习、对比学习、预训练--微调与迁移/元学习的目标；
  \item likelihood、normalizing flow、VAE ELBO、GAN、DDPM ELBO 与 score matching；
  \item Bellman equation、policy iteration、actor-critic 与 policy gradient theorem。
\end{itemize}

\DeepTextbooks

\subsection{人工智能中的优化方法}

官方大纲要求四组内容：

\begin{enumerate}
  \item 凸集与凸函数；
  \item 次微分及其性质、凸函数的重要性质；
  \item 梯度下降、随机梯度下降与随机坐标下降；
  \item 动量加速与自适应学习率。
\end{enumerate}

指定教材为 \textcite{nesterov2018convex,beck2017first}。公开卷中常把这些对象组合为 composite optimization：

\[
  \min_{x\in \R^d} F(x) = f(x) + R(x),
  \qquad
  x_{k+1} = \prox_{\gamma R}\bigl(x_k-\gamma g(x_k,\xi_k)\bigr).
\]

必须掌握以下链条：

\[
  \text{强凸/光滑}
  \Longrightarrow
  \text{梯度映射收缩}
  \Longrightarrow
  \text{prox 非扩张}
  \Longrightarrow
  \text{一阶递推}
  \Longrightarrow
  \text{噪声地板与复杂度}.
\]

\OptimizationTextbooks

\subsection{自然语言处理}

官方大纲的六项内容是：

\begin{enumerate}
  \item 自然语言处理任务的统计建模假设与方法；
  \item Word2vec 的数学原理；
  \item Transformer 中 self-attention 的原理与实现；
  \item 大模型预训练对齐方法的原理、局限与改进；
  \item Scaling Law 的基本原理；
  \item 基于 embedding 的知识库推理方法、局限与改进。
\end{enumerate}

指定教材为 \textcite{jurafsky2026slp,goldberg2017nnnlp}。公开卷在此基础上还出现 RAG、MLN、topic model、contrastive learning、长上下文 Transformer 和 text-to-image diffusion system design；这些属于考纲主线的具体应用和延伸。

复习时把 NLP 题分成三类：

\begin{itemize}
  \item 概念题：n-gram smoothing、representation learning、top-\(p\)、Transformer 优势；
  \item 数学题：InfoNCE 梯度、scaling law、MLN 概率、supervised LDA joint；
  \item 系统题：RAG、RLHF、长上下文、信息抽取全局一致性、文生图 diffusion。
\end{itemize}

\NLPTextbooks

\printbibliography[heading=bibintoc,title={参考文献与官方来源}]

\end{document}
